\documentclass[10pt,a4paper]{amsart}
\usepackage{amsmath,amssymb,amsthm,graphicx}
% Retained minimal preamble for the Stacks Project chapter flat.tex (More on Flatness).
% Verbatim excerpt of the chapter preamble: Xy-pic setup (native preamble lines 14-18) and the native
% theorem declarations with the mathematical macros (native preamble lines 154-217). Copied mechanically;
% nothing added, removed or redefined. The wrapper class and the amsmath/amssymb/amsthm/graphicx base
% packages are supplied by the assembler, so the class/externaldocument part of the native preamble is omitted.
% For commutative diagrams we use Xy-pic
\usepackage[all]{xy}

% We use 2cell for 2-commutative diagrams.
\xyoption{2cell}
% Theorem environments.
%
\theoremstyle{plain}
\newtheorem{theorem}[subsection]{Theorem}
\newtheorem{proposition}[subsection]{Proposition}
\newtheorem{lemma}[subsection]{Lemma}

\theoremstyle{definition}
\newtheorem{definition}[subsection]{Definition}
\newtheorem{example}[subsection]{Example}
\newtheorem{exercise}[subsection]{Exercise}
\newtheorem{situation}[subsection]{Situation}

\theoremstyle{remark}
\newtheorem{remark}[subsection]{Remark}
\newtheorem{remarks}[subsection]{Remarks}

\numberwithin{equation}{subsection}

% Macros
%
\def\lim{\mathop{\mathrm{lim}}\nolimits}
\def\colim{\mathop{\mathrm{colim}}\nolimits}
\def\Spec{\mathop{\mathrm{Spec}}}
\def\Hom{\mathop{\mathrm{Hom}}\nolimits}
\def\Ext{\mathop{\mathrm{Ext}}\nolimits}
\def\SheafHom{\mathop{\mathcal{H}\!\mathit{om}}\nolimits}
\def\SheafExt{\mathop{\mathcal{E}\!\mathit{xt}}\nolimits}
\def\Sch{\mathit{Sch}}
\def\Mor{\mathop{\mathrm{Mor}}\nolimits}
\def\Ob{\mathop{\mathrm{Ob}}\nolimits}
\def\Sh{\mathop{\mathit{Sh}}\nolimits}
\def\NL{\mathop{N\!L}\nolimits}
\def\CH{\mathop{\mathrm{CH}}\nolimits}
\def\proetale{{pro\text{-}\acute{e}tale}}
\def\etale{{\acute{e}tale}}
\def\QCoh{\mathit{QCoh}}
\def\Ker{\mathop{\mathrm{Ker}}}
\def\Im{\mathop{\mathrm{Im}}}
\def\Coker{\mathop{\mathrm{Coker}}}
\def\Coim{\mathop{\mathrm{Coim}}}

% Boxtimes
%
\DeclareMathSymbol{\boxtimes}{\mathbin}{AMSa}{"02}

%
% Macros for moduli stacks/spaces
%
\def\QCohstack{\mathcal{QC}\!\mathit{oh}}
\def\Cohstack{\mathcal{C}\!\mathit{oh}}
\def\Spacesstack{\mathcal{S}\!\mathit{paces}}
\def\Quotfunctor{\mathrm{Quot}}
\def\Hilbfunctor{\mathrm{Hilb}}
\def\Curvesstack{\mathcal{C}\!\mathit{urves}}
\def\Polarizedstack{\mathcal{P}\!\mathit{olarized}}
\def\Complexesstack{\mathcal{C}\!\mathit{omplexes}}
% \Pic is the operator that assigns to X its picard group, usage \Pic(X)
% \Picardstack_{X/B} denotes the Picard stack of X over B
% \Picardfunctor_{X/B} denotes the Picard functor of X over B
\def\Pic{\mathop{\mathrm{Pic}}\nolimits}
\def\Picardstack{\mathcal{P}\!\mathit{ic}}
\def\Picardfunctor{\mathrm{Pic}}
\def\Deformationcategory{\mathcal{D}\!\mathit{ef}}

% --- appended 2026-10-12 (Chief): the project own definitions of the reference/slogan/history
% display environments, copied mechanically from the shared Stacks preamble lines 8-12
% (\usepackage{verbatim} plus the three \newenvironment lines, which the chapter preamble keeps
% verbatim from the project preamble). Nothing added, removed or redefined: these environments
% carry no mathematical content and the project itself typesets them as comments. Without these
% three lines a counted statement whose own text uses them fails with Environment slogan undefined.
\usepackage{verbatim}
\newenvironment{reference}{\comment}{\endcomment}
\newenvironment{slogan}{\comment}{\endcomment}
\newenvironment{history}{\comment}{\endcomment}
\usepackage{cleveref}
\crefname{originalreference}{original reference}{original references}
\Crefname{originalreference}{Original reference}{Original references}
\makeatletter
\newlabel{situation-mod-injective}{{0AT0}{source}{}{sourcepointer-1}{}}
\newlabel{situation-mod-injective@cref}{{[originalreference][0][]0AT0}{[1][0][]source}{}{}{}}
\makeatother
\crefname{originallemma}{original lemma}{original lemmas}
\Crefname{originallemma}{Original lemma}{Original lemmas}
\makeatletter
\newlabel{algebra-lemma-finitely-presented-over-subring}{{0561}{source}{}{sourcepointer-2}{}}
\newlabel{algebra-lemma-finitely-presented-over-subring@cref}{{[originallemma][0][]0561}{[1][0][]source}{}{}{}}
\makeatother
\title{The injectivity situation reduces to localizations of polynomial algebras with finitely presented modules}
\author{The Stacks Project Authors}
\date{stacks-project:flat.tex:local-snapshot-20261011; GNU FDL 1.2 or later}
\begin{document}
\maketitle
\noindent\textit{Editorial note.} All delivered mathematical text is the author's own, copied from the pinned Stacks Project chapter file flat.tex; no statement or proof step was rewritten, reordered or removed.\par\medskip
\noindent\textit{Editorial note.} The preamble is the chapter's own retained minimal preamble (sources/stacks-flat/retained-preamble.tex), a verbatim excerpt of the native Xy-pic setup and of the native theorem declarations and mathematical macros, followed by the project own definitions of the reference, slogan and history display environments (verbatim package plus three newenvironment lines) that the retained file carries at its end and that the project itself typesets as comments. No macro was invented or redefined, and the wrapper class and the amsmath/amssymb/amsthm/graphicx base packages are supplied by the assembler.\par\medskip
\noindent\textit{Editorial note.} The counted statement and proof contain no cross reference, no citation and no included file, so the delivered body is self contained.\par\medskip
\noindent\textit{Editorial note.} Canonical location of the counted proposition: Stacks tag 0AT2 (sources/stacks-flat/pinned-original-source.tex line 7045).\par\medskip
\noindent\textit{Allowed standard prerequisites (editorial).}\begin{itemize}
\item Finite presentation over a subring from finite presentation over a quotient. Algebra, Lemma labelled algebra-lemma-finitely-presented-over-subring, cited at proof lines 7057-7059. Lets one replace B by the localization of a polynomial algebra.
\item The situation and its conclusion. Same chapter, Situation labelled situation-mod-injective, invoked in the statement. Fixes the objects and the conclusion being reduced.
\item Noetherianity of the special fibre. The fact that B/m\_A B is essentially of finite type over a field, used at proof lines 7062-7066. Allows a finite choice of generators of K cutting out N modulo the maximal ideal.
\end{itemize}\medskip
\ifdefined\hypertarget\hypertarget{sourcepointer-1}{}\fi
\noindent\textit{Source reference.} Original reference 0AT0: sources/stacks-flat/pinned-original-source.tex, source lines 7017--7017. Same-chapter pointer into the author environment situation at source line 7017; invoked at proof line 7054 of the counted proof. It is retained only as a navigation pointer; every load-bearing step is inside the delivered span.\par\medskip
\ifdefined\hypertarget\hypertarget{sourcepointer-2}{}\fi
\noindent\textit{Source reference.} Original lemma 0561: sources/stacks-chow/tags.txt, source lines 6694--6694. Cross-chapter standard prerequisite located in the project tag list staged in the local snapshot; invoked at proof line 7059 of the counted proof. It is retained only as a navigation pointer; every load-bearing step is inside the delivered span.\par\medskip
\begin{lemma}
\label{lemma-reduce-finite-type-injective-into-flat-mod-m}
Let $A_0$ be a local ring. If the lemma holds for every
Situation \ref{situation-mod-injective} with $A = A_0$, with $B$ a
localization of a polynomial algebra over $A$, and $N$ of finite presentation
over $B$, then the lemma holds for every
Situation \ref{situation-mod-injective} with $A = A_0$.
\end{lemma}

\begin{proof}
Let $A \to B$, $u : N \to M$ be as in Situation \ref{situation-mod-injective}.
Write $B = C/I$ where $C$ is the localization of a polynomial algebra
over $A$ at a prime. If we can show that $N$ is finitely presented as
a $C$-module, then a fortiori this shows that $N$ is finitely presented
as a $B$-module (see
Algebra, Lemma \ref{algebra-lemma-finitely-presented-over-subring}).
Hence we may assume that $B$ is the localization of a polynomial algebra.
Next, write $N = B^{\oplus n}/K$ for some submodule $K \subset B^{\oplus n}$.
Since $B/\mathfrak m_AB$ is Noetherian (as it is essentially of finite type
over a field), there exist finitely many elements $k_1, \ldots, k_s \in K$
such that for $K' = \sum Bk_i$ and $N' = B^{\oplus n}/K'$ the
canonical surjection $N' \to N$ induces an isomorphism
$N'/\mathfrak m_AN' \cong N/\mathfrak m_AN$.
Now, if the lemma holds for the composition $u' : N' \to M$,
then $u'$ is injective, hence $N' = N$ and $u' = u$. Thus the lemma holds for
the original situation.
\end{proof}

\end{document}
